\begin{Exercise}[label={ex:caratteri}]
  Scrivere un programma che prenda in ingresso un carattere (e.g., \texttt{a}, \texttt{b}, \texttt{c}) ed un numero intero positivo $n$. Il programma dovrà stampare a video il carattere inserito, formattato come carattere e come intero. Poi dovrà stampare a video il carattere $n$ posizioni successive a quello inserito. Ad esempio, se il carattere inserito è la lettera \texttt{a} ed $n = 3$, il programma dovrà stampare a video il carattere \texttt{d}.
\end{Exercise}

\begin{Answer}[ref={ex:caratteri}]
  \lstinputlisting[style=cstyle]{src/caratteri.c}
\end{Answer}